% Spanish – Geography, Ser/Estar, Weather & Description (Study Guide)
% Everything from chat: exercises 5–7, 10–13, 15, 17–18, 20–25, memory hooks.
% Compile: pdflatex geography-weather.tex

\documentclass[9pt,a4paper]{article}
\usepackage[utf8]{inputenc}
\usepackage[T1]{fontenc}
\usepackage[margin=1cm]{geometry}
\usepackage{multicol}
\usepackage{enumitem}
\setlist{nosep,leftmargin=*}
\usepackage{booktabs}
\usepackage{parskip}
\usepackage{xcolor}
\usepackage{fancyhdr}
\usepackage{microtype}

\definecolor{spanishgreen}{RGB}{0,100,80}
\definecolor{charcoal}{RGB}{51,51,51}
\definecolor{notegrey}{RGB}{90,90,90}

\newcommand{\tip}[1]{\textcolor{spanishgreen}{\textbf{#1}}}
\newcommand{\wrong}[1]{\textcolor{red}{#1}}
\setlength{\parskip}{0.25em}
\pagestyle{fancy}
\fancyhf{}
\fancyfoot[L]{\small\color{notegrey} Spanish for Beginners}
\fancyfoot[R]{\small\color{notegrey}\thepage}
\raggedright

\begin{document}
\begin{center}
\vspace{0.2em}
{\large\bfseries\color{spanishgreen} Study Guide: Everything From Chat}\\[0.15em]
\footnotesize\textcolor{notegrey}{Ser/estar, hay, questions, muy/mucho, weather, seasons, geography, adjectives, ex. 5--25.}\\[0.3em]
\end{center}

\begin{multicols}{2}

\section*{\color{spanishgreen}Ser vs Estar}
\textbf{When:} Ser = what/it is (identity). Estar = where/how (location, state).
\textbf{Saying:} ``Where it is and how it is $\to$ \textbf{estar}. What it is $\to$ \textbf{ser}.''
\textbf{Mnemonics:} Ser = DOOT (Description, Occupation, Origin, Time). Estar = PLACE (Position, Location, Action, Condition, Emotion).

\subsection*{Conjugations}
\footnotesize
Ser: soy, eres, es, somos, sois, son. \quad Estar: estoy, est\'as, est\'a, estamos, est\'ais, est\'an.
\textbf{Accent:} est\'a (is) vs esta (this).

\subsection*{Ex. 5 -- Ser/Estar choice}
Location $\to$ \textbf{est\'a}: M\'exico est\'a en Am\'erica del Norte. Puerto Rico est\'a en el Caribe. Barcelona est\'a en la costa. \\
Identity $\to$ \textbf{es}: Valpara\'iso es una ciudad muy tur\'istica. Mallorca es una isla. Lima es la capital de Per\'u.

\subsection*{Ex. 6 -- Complete logically}
(d) \textbf{Est\'a} + place: \textit{La ciudad donde vivo est\'a en el medioeste (de EE.UU.).} \\
(e) \textbf{Es} + description: \textit{La ciudad donde vivo es grande / muy tur\'istica.} \\
Wrong: ``est\'a es la capital'' or ``es al norte'' (swap: location=estar, description=ser).

\section*{\color{spanishgreen}Hay, es/son, est\'a/est\'an (Ex. 7)}
\textbf{Hay} = there is/are. \textbf{Ser} = identity. \textbf{Estar} = location. \\
Ejemplos: En Guatemala \textbf{hay} ruinas mayas. Quito \textbf{es} la capital de Ecuador. Machu Picchu \textbf{est\'a} en Per\'u. La moneda de Bolivia \textbf{es} el boliviano. En Venezuela \textbf{hay} petr\'oleo.

\section*{\color{spanishgreen}Question words (Ex. 11, 12)}
\begin{center}\footnotesize
\begin{tabular}{@{}ll@{}}
\toprule
\textbf{Word} & \textbf{Use when answer is\ldots} \\
\midrule
¿D\'onde? & a place \\
¿C\'omo? & description, ``what's it like'' \\
¿Qu\'e? & definition, ``what is X'' \\
¿Cu\'al/Cu\'ales? & which one(s), name/choice \\
¿Cu\'anto(s)/Cu\'anta(s)? & number, amount \\
\bottomrule
\end{tabular}
\end{center}
\textbf{Saying:} C\'omo = like/how (description). Qu\'e = what (the thing/definition). \\
\textbf{Ex. 11 Bolivia:} ¿D\'onde est\'a Bolivia? (place). ¿C\'omo es el clima? (description). ¿Qu\'e es el ajiaco? (definition). ¿Cu\'al es la bebida t\'ipica? (name). ¿Cu\'ales son las ciudades? (which). ¿Cu\'antas lenguas? (number). \\
\textbf{Ex. 12a:} Qu\'e = bachata, flamenco (definition). Cu\'al = pa\'is m\'as peque\~no, r\'io m\'as largo. Cu\'ales = playas C\'adiz, productos Cuba. \\
\textbf{Ex. 12b answers:} Bachata = m\'usica/baile Rep. Dominicana. Playas C\'adiz = La Caleta, La Victoria, Cortadura. Productos Cuba = tabaco, ron. Flamenco = arte Andaluc\'ia (cante, toque, baile). Pa\'is m\'as peque\~no Am\'erica Central = El Salvador. R\'io m\'as largo Am\'erica del Sur = Amazonas.

\section*{\color{spanishgreen}Muy vs mucho/a/os/as (Ex. 10, 18)}
\begin{itemize}
\item \textbf{Muy} = very. Goes with \textbf{adjectives}: muy antiguo, muy importante, muy poco.
\item \textbf{Mucho/mucha/muchos/muchas} = much/many. Agree with \textbf{noun}: muchos lugares, muchas playas, mucho calor.
\item \textbf{Adverb} (a lot): \textit{llueve mucho}, \textit{no llueve mucho}.
\end{itemize}
Before adjective $\to$ muy. Before noun (or ``rains a lot'') $\to$ mucho/muchos/muchas. \\
\textbf{Ex. 10:} muy antigua, \textbf{muchos} lugares (masc!), muchas playas, muchos edificios, muy importante, muy simp\'atica, muchos templos, muy conocido. \\
\textbf{Ex. 18 errors:} \textbf{muchos} climas (not muchas). \textbf{muy} secos (not muchos). llueve \textbf{mucho} (not muy). hace \textbf{mucho} fr\'io/calor (not muy). \textbf{muchas} zonas. \textbf{muy} poco. Last blank before mundo = \textbf{del}.

\section*{\color{spanishgreen}Superlatives (Ex. 13)}
Pattern: \textbf{el/la/los/las} + noun + \textbf{m\'as} + adjective + \textbf{de/del}.
\begin{itemize}
\item \textit{El Nilo es el r\'io m\'as largo de \'Africa.}
\item \textit{Rusia es el pa\'is m\'as grande del mundo.}
\item Before \textit{mundo}: \textbf{del} (de + el). Before place name (no article): \textbf{de} (de Europa, de Latinoam\'erica).
\end{itemize}
Middle blank is always \textbf{m\'as}. Before mundo = \textbf{del}. Spell \textbf{m\'as} with accent (mas = but). \\
\textbf{Ex. 13 full:} (a) el, m\'as, de (b) el, m\'as, del (c) el, m\'as, de (d) el, m\'as, del (e) las, m\'as, del (f) la, m\'as, del (g) la, m\'as, del (h) El, m\'as, del (i) los, m\'as, de (j) las, m\'as, del. Pa\'is/isla = el/la. Monta\~nas/pir\'amides = las.

\section*{\color{spanishgreen}Weather (Ex. 15, 17)}
\begin{itemize}
\item \textbf{Hace} + noun: hace calor, hace fr\'io, hace sol, hace viento, hace buen/mal tiempo.
\item \textbf{Estar} + adjective: est\'a nublado, est\'a soleado.
\item \textbf{Hay} + noun: hay nubes (not ``est\'a nubes'').
\item \textbf{Verbs}: llueve, nieva.
\end{itemize}
\textbf{Ex. 15:} ``Est\'a nubes'' wrong $\to$ \textbf{Hay nubes} (nubes = noun). Hace calor/fr\'io/viento. Est\'a nublado. Es tropical/seco/h\'umedo/c\'alido. \\
\textbf{Ex. 17 Spain:} Santiago llueve. Bilbao est\'a nublado. Huesca nieva. Soria hace fr\'io. Mallorca hace viento. Sevilla hace sol. C\'adiz hace calor. (Match to map.)

\section*{\color{spanishgreen}Seasons \& vocab (Ex. 20)}
\begin{itemize}
\item \textbf{Invierno}: hace fr\'io, nieve, nieva, abrigo, diciembre.
\item \textbf{Verano}: hace calor, hace sol, playa, vacaciones, agosto.
\item \textbf{Primavera}: flores, hace buen tiempo, llueve, abril, mayo.
\item \textbf{Oto\~no}: hojas, hace fresco, viento, lluvia, octubre.
\end{itemize}
\textbf{English + memory:} Invierno = \textit{in} the cold. Verano = \textit{ver} (see) the sun. Primavera = \textit{first}/prime. Oto\~no = \textit{O}ctober, orange leaves.

\subsection*{Ex. 18 -- Clima en Espa\~na (correct)}
Espa\~na es un pa\'is con \textbf{muchos} climas diferentes. Los veranos son \textbf{muy} secos, no llueve \textbf{mucho}, no hace \textbf{mucho} fr\'io. En el norte llueve \textbf{mucho}. En el interior hace \textbf{mucho} calor y \textbf{mucho} fr\'io. En \textbf{muchas} zonas del sur llueve \textbf{muy} poco, hace \textbf{mucho} calor.

\section*{\color{spanishgreen}Geography nouns (Ex. 21)}
\begin{itemize}
\item r\'io (river), isla (island), monta\~na (mountain), capital (capital), cordillera (mountain range), bebida (drink).
\item \textbf{Ex. 21:} (a) r\'io (b) isla (c) monta\~na (d) ciudad (e) cordillera (f) capital (g) bebida.
\end{itemize}

\section*{\color{spanishgreen}Adjectives with nouns (Ex. 22)}
\begin{itemize}
\item \textbf{Clima}: templado, h\'umedo, fr\'io, seco, tropical, bueno, interesante, importante.
\item \textbf{Monta\~na}: alta, grande, bonita, famosa, tur\'istica, tropical (feminine: alta, bonita).
\item \textbf{Pa\'is}: grande, poblado, importante, bonito, famoso, tur\'istico, tropical.
\item \textbf{Comida}: buena, t\'ipica, famosa, interesante, tropical, tur\'istica (feminine).
\end{itemize}
\textbf{Ex. 22:} Unclick \textbf{t\'ipica} for MONTA\~NA (wrong). Click for CLIMA: templado, h\'umedo, fr\'io, seco, tropical, bueno, interesante, importante, grande. MONTA\~NA: alta, grande, bonita, preciosa, famosa, importante, interesante, tur\'istica, tropical (+ fr\'ia, h\'umeda, seca, templada, poblada). PA\'IS: all except alto. COMIDA: buena, t\'ipica, famosa, preciosa, interesante, templada, tropical, tur\'istica, seca, grande. Aim 26/26.

\section*{\color{spanishgreen}Expressions (Ex. 24)}
\begin{itemize}
\item un pa\'is: muy poblado / con pocos habitantes
\item un r\'io: muy largo / caudaloso
\item una lengua: oficial / materna
\item una monta\~na: muy alta / nevada
\item un plato: t\'ipico / tradicional
\item una isla: tropical / con playas
\item un desierto: muy grande / \'arido
\item un clima: templado / tropical
\item una ciudad: muy tur\'istica / antigua
\item un palacio: antiguo / hist\'orico
\end{itemize}

\section*{\color{spanishgreen}Ex. 25 -- Chat Leda/Ana (V o F)}
1. Do\~nana en Madrid? \textbf{F}. (Andaluc\'ia.) \quad 2. Retiro en Madrid? \textbf{V}. \quad 3. Acueducto en Segovia? \textbf{V}. \quad 4. Giralda en Granada? \textbf{F}. (Sevilla. Alhambra = Granada.) \quad 5. Playas en Almer\'ia? \textbf{V}.

\section*{\tip{Quick recap}}
\begin{itemize}
\item Ser = what/identity. Estar = where/state. Hay = there is/are.
\item Question: place $\to$ D\'onde; description $\to$ C\'omo; definition $\to$ Qu\'e; which $\to$ Cu\'al; number $\to$ Cu\'anto(s).
\item Muy + adjective. Mucho(s)/mucha(s) + noun; llueve mucho.
\item Superlative: el/la/los/las + m\'as + adj + de/del.
\item Weather: hace calor/fr\'io/sol/viento; est\'a nublado; hay nubes; llueve, nieva.
\item Ex. 25: Do\~nana $\to$ Andaluc\'ia. Giralda $\to$ Sevilla. Alhambra $\to$ Granada.
\end{itemize}

\end{multicols}
\end{document}
